% v2026a baseline; generated by R/44_Ideology_Diagnosis.R; external sources: data/external/SOURCES.csv
\begin{table}[ht]
\caption{Correlation of the LLM Measures with External Indicators}
\label{tab:external}
\footnotesize
\centering
\resizebox{\ifdim\width>\linewidth\linewidth\else\width\fi}{!}{%
\begin{tabular}{lrrrr}
\toprule
 & \multicolumn{2}{c}{Populism (LLM)} & \multicolumn{2}{c}{Ideology (LLM)} \\
\cmidrule(lr){2-3} \cmidrule(lr){4-5}
External indicator & $r$ & N & $r$ & N \\
\midrule
\multicolumn{5}{l}{\emph{CMP (Manifesto Project 2026a)}} \\
\quad rile & -0.12 & 3,349 & 0.65 & 3,209 \\
\quad rile, logit scale & -0.14 & 3,342 & 0.64 & 3,202 \\
\multicolumn{5}{l}{\emph{CHES 1999--2024}} \\
\quad lrgen (general left-right) & -0.14 & 875 & 0.85 & 847 \\
\quad lrecon (economic left-right) & -0.43 & 875 & 0.69 & 847 \\
\quad galtan (GAL-TAN) & 0.17 & 875 & 0.72 & 847 \\
\multicolumn{5}{l}{\emph{ParlGov}} \\
\quad left\_right & -0.36 & 2,010 & 0.82 & 1,916 \\
\multicolumn{5}{l}{\emph{POPPA 2018}} \\
\quad populism & 0.66 & 177 & 0.29 & 173 \\
\quad left-right, overall & -0.14 & 189 & 0.84 & 185 \\
\quad left-right, economic & -0.46 & 185 & 0.67 & 181 \\
\multicolumn{5}{l}{\emph{Global Party Survey 2019}} \\
\quad V8 populist rhetoric & 0.39 & 276 & 0.45 & 273 \\
\quad V4 economic left-right & -0.39 & 275 & 0.64 & 272 \\
\quad V6 social conservatism & -0.02 & 276 & 0.69 & 273 \\
\multicolumn{5}{l}{\emph{RFPOPI (Celico et al. 2024)}} \\
\quad populism, rhetorical & 0.24 & 1,852 & 0.22 & 1,769 \\
\quad populism, ideational & 0.63 & 1,852 & 0.15 & 1,769 \\
\multicolumn{5}{l}{\emph{PopuList 4.0}} \\
\quad populist vs.\ not, AUC & 0.80 & 1,582 & & \\
\quad far right vs.\ far left, AUC & & & 0.99 & 384 \\
\bottomrule
\end{tabular}}
\par\smallskip
\begin{minipage}{\linewidth}\scriptsize Pearson correlations at the manifesto level; N = manifestos with both scores. Matching: the right-left index (rile) of the Comparative Manifesto Project (CMP) and its logit version (Lowe et al.\ 2011, from the category counts of the Manifesto Project 2026a) from the same manifesto; Chapel Hill Expert Survey (CHES) trend file for the same party in the wave nearest to the election year, at most two years away (lrgen: general left-right; lrecon: economic left-right; galtan: green-alternative-libertarian (GAL) vs.\ traditional-authoritarian-nationalist (TAN)); ParlGov left\_right as the party's time-invariant mean; Populism and Political Parties Expert Survey (POPPA) 2018 (Meijers and Zaslove) and Global Party Survey (GPS) 2019 (Norris): each party's election nearest to the survey year, at most two years away, party identifiers linked to CMP codes via Partyfacts; PopuList 4.0 matched via Partyfacts and classified for the election year. AUC: probability that a manifesto of the first group scores higher than one of the second; populist vs.\ not uses all manifestos with a populism score from PopuList's countries since 1989, parties not on the list counted as non-populist (295 / 1,287); far right vs.\ far left: 193 / 191. RFPOPI (random-forest populism indicators of Celico, Rode and Rodr\'{\i}guez Carre\~{n}o 2024): rhetorical and ideational populism of the party in the same election year (or within one year), linked via Partyfacts.  Details of the matching: Online Appendix, Table~\ref{tab:oa1b}.\end{minipage}
\end{table}
